\section{Motivation}
\label{sec:motivation}

\subsection{Lifetimes heterogeneity within one prompt}
The coding example in Figure~\ref{fig:lowering} illustrates why treating an entire prompt as one reuse unit loses useful information. Consider an iterative implementation task in which \texttt{spec} and \texttt{plan} remain fixed, \texttt{advice} is replaced after a failed test, and \texttt{attempts} accumulates earlier implementations and test results. These fields belong to the same invocation, yet their next values have different relationships to the current ones. The specification and plan can be copied, part of the attempt history remains known, and the next feedback is unavailable until testing completes. Framework instructions may persist even across sessions, whereas the specification is private to one task.

The fact-checking example in Figure~\ref{fig:heterogeneity-sample} has the same structure: the claim persists across a verification session, evidence accumulates, and a new sub-question or search result can replace the corresponding input. Moreover, the same claim or evidence may be passed between call sites. This creates two distinct dimensions of heterogeneity: how a binding evolves over time, and where its value is shared. A field can be stable within one session without being a program-wide constant, and it can be shared across sites without remaining unchanged forever. These properties must be inferred from values and execution history; a parameter name or type alone does not establish them.

\subsection{Declaration order hides reusable content}
Frameworks organize fields to express an operation's interface. That order need not expose the longest reusable token prefix. Suppose the coding request is rendered as
\[
P_t = H \Vert S \Vert L \Vert A_t \Vert E_t \Vert R,
\]
where $H$ contains fixed leading instructions and metadata, $S$ is the specification, $L$ is the plan, $A_t$ is current advice, $E_t$ is the attempt history, and $R$ contains trailing instructions. Each field includes its label and delimiters. When advice changes, prefix matching stops within $A_t$ even if most of $E_t$ remains identical. Keeping the entire previous request in GPU memory cannot recover this lost reuse: the old history tensors were computed under the old advice.

Placing the attempt history before advice instead gives
\[
P'_t = H \Vert S \Vert L \Vert E_t \Vert A_t \Vert R.
\]
If consecutive serialized histories share an initial segment, the reusable prefix can now extend through that segment. Appending a list element may replace a closing delimiter, so the reusable region is the actual common token prefix, not necessarily the complete previous field. This example requires no prediction of the next advice. It exploits the part of the next request that is already known.

Cross-call-site sharing introduces another boundary. Two sites may receive identical leading bindings but have different headers. A site-specific header before those bindings prevents sharing across sites at that point. Moving shared bindings ahead of that header can expose further reuse when the system prompts, serialization, and preceding tokens also match. Identical binding values alone are insufficient. These observations motivate learning a consistent layout that considers both sharing scope and evolution, while retaining labels and field contents. 

\subsection{Reuse opportunities need timely residency}
A longer common prefix is useful only if its KV state is available. Recency alone cannot distinguish an intermediate result that another call will soon consume from a fresh diagnostic that the next attempt will replace. Nor does a recent access imply continued usefulness after its session terminates. Conversely, a prefix can remain useful throughout a tool execution despite receiving no accesses during that interval. Continuum addresses retention across such pauses through cache time-to-live, while CacheScout learns execution transitions to guide retention and prefetching~\cite{li2025continuum,zhang2026cachescout}.

Binding heterogeneity supplies information beyond the identity of a future call. Predicting that \texttt{implement} will run again doesn't utilize its full potential. Its specification and plan can be reconstructed from the live session, the history contributes an existing prefix, and new advice remains unresolved. Re-layout can place these reconstructible inputs before the first unknown field, increasing the amount of KV that can be prepared before arrival. Call-transition and timing observations could determine which candidate prefixes are likely to be needed and when promotion should begin.
